\documentclass[openany]{book}

\usepackage[margin=1in]{geometry}
\usepackage{amsmath,amsfonts,amsthm, amssymb}
\usepackage{yhmath}
\usepackage{mathrsfs}
\usepackage{mathtools}
\usepackage{xcolor}
\usepackage{graphicx}
\usepackage{comment}
\usepackage{tikz-cd}
\usepackage{quiver}
\renewcommand{\familydefault}{ppl}
\newcommand{\tr}{\text{tr}}
\newcommand{\R}{\mathbb{R}}
\newcommand{\E}{\mathbb{E}}
\newcommand{\Z}{\mathbb{Z}}
\newcommand{\C}{\mathbb{C}}
\newcommand{\F}{\mathbb{F}}
\newcommand{\la}{\langle}
\newcommand{\ra}{\rangle}
\newcommand{\colim}{\text{colim}}
\DeclareMathOperator{\im}{im}
\let\oldemptyset\emptyset
\let\emptyset\varnothing
\newcommand{\tor}{\text{Tor}}
\newcommand{\id}{\text{id}}
\newcommand{\ext}{\text{Ext}}
\newcommand{\ptop}{\text{PTop}}
\newcommand{\pt}{\text{pt}}
\newcommand{\ach}{\text{Ach}}
\newcommand{\Q}{\mathbb{Q}}
\newcommand{\gal}{\text{Gal}}


\usepackage{thmtools,thm-restate}

% Fixing mdframed skip below
% See https://tex.stackexchange.com/a/292090/143086
\usepackage[framemethod=TikZ]{mdframed}
\usepackage{xpatch}
\makeatletter
\xpatchcmd{\endmdframed}
	{\aftergroup\endmdf@trivlist\color@endgroup}
	{\endmdf@trivlist\color@endgroup\@doendpe}
	{}{}
\makeatother

\definecolor{huilightpink}{HTML}{fff2fe}
\definecolor{huidarkpink}{HTML}{d955b7}
\declaretheoremstyle[
	mdframed={
		backgroundcolor=huilightpink,
		linecolor=huidarkpink,
		rightline=false,
		topline=false,
		bottomline=false,
		linewidth=2pt,
		innertopmargin=5pt,
		innerbottommargin=8pt,
		innerleftmargin=8pt,
		leftmargin=-2pt,
		skipbelow=2pt,
		nobreak
	},
	headfont=\normalfont\bfseries\color{huidarkpink}
]{huipinkbox}
\declaretheorem[style=huipinkbox,name=Theorem,within=chapter]{thm}
\declaretheorem[style=huipinkbox,name=Theorem,sibling=thm]{theorem}




\begin{comment}
\definecolor{huilightyellow}{HTML}{fff5d6}
\definecolor{huidarkyellow}{HTML}{fcad03}
\declaretheoremstyle[
	mdframed={
		backgroundcolor=huilightyellow,
		linecolor=huidarkyellow,
		rightline=false,
		topline=false,
		bottomline=false,
		linewidth=2pt,
		innertopmargin=5pt,
		innerbottommargin=8pt,
		innerleftmargin=8pt,
		leftmargin=-2pt,
		skipbelow=2pt,
		nobreak
	},
	headfont=\normalfont\bfseries\color{huidarkyellow}
]{huiyellowbox}
\declaretheorem[style=huiyellowbox,name=Proposition,within=chapter]{prop}
\end{comment}



\definecolor{huilightpurple}{HTML}{faf2ff}
\definecolor{huidarkpurple}{HTML}{912ed9}
\declaretheoremstyle[
	mdframed={
		backgroundcolor=huilightpurple,
		linecolor=huidarkpurple,
		rightline=false,
		topline=false,
		bottomline=false,
		linewidth=2pt,
		innertopmargin=5pt,
		innerbottommargin=8pt,
		innerleftmargin=8pt,
		leftmargin=-2pt,
		skipbelow=2pt,
		nobreak
	},
	headfont=\normalfont\bfseries\color{huidarkpurple}
]{huipurplebox}
\declaretheorem[style=huipurplebox,name=Proposition,within=chapter]{prop}



% \definecolor{huilightpurple}{HTML}{faf2ff}
% \definecolor{huidarkpurple}{HTML}{912ed9}
% \declaretheoremstyle[
% 	mdframed={
% 		backgroundcolor=huilightpurple,
% 		linecolor=huidarkpurple,
% 		rightline=false,
% 		topline=false,
% 		bottomline=false,
% 		linewidth=2pt,
% 		innertopmargin=5pt,
% 		innerbottommargin=8pt,
% 		innerleftmargin=8pt,
% 		leftmargin=-2pt,
% 		skipbelow=2pt,
% 		nobreak
% 	},
% 	headfont=\normalfont\bfseries\color{huidarkpurple}
% ]{huipurplebox}
\declaretheorem[style=huipurplebox,name=Lemma,within=chapter]{lem}


\definecolor{lightpink}{HTML}{f0f6fc}
\definecolor{darkpink}{HTML}{2c72b8}
\declaretheoremstyle[
	mdframed={
		backgroundcolor=lightpink,
		linecolor=darkpink,
		rightline=false,
		topline=false,
		bottomline=false,
		linewidth=2pt,
		innertopmargin=5pt,
		innerbottommargin=8pt,
		innerleftmargin=8pt,
		leftmargin=-2pt,
		skipbelow=2pt,
		nobreak
	},
	headfont=\normalfont\bfseries\color{darkpink}
]{pinkbox}
\declaretheorem[style=pinkbox,name=Definition,within=chapter]{defn}


\definecolor{huilightblue}{HTML}{edf9ff}
\definecolor{huidarkblue}{HTML}{4b79db}
\declaretheoremstyle[
	mdframed={
		backgroundcolor=huilightblue,
		linecolor=huidarkblue,
		rightline=false,
		topline=false,
		bottomline=false,
		linewidth=2pt,
		innertopmargin=5pt,
		innerbottommargin=8pt,
		innerleftmargin=8pt,
		leftmargin=-2pt,
		skipbelow=2pt,
		nobreak
	},
	headfont=\normalfont\bfseries\color{huidarkblue}
]{huiblueblox}
\declaretheorem[style=huiblueblox,name=Example,within=chapter]{example}



% \definecolor{huilightblue}{HTML}{edf9ff}
% \definecolor{huidarkblue}{HTML}{4b79db}
% \declaretheoremstyle[
% 	mdframed={
% 		backgroundcolor=huilightblue,
% 		linecolor=huidarkblue,
% 		rightline=false,
% 		topline=false,
% 		bottomline=false,
% 		linewidth=2pt,
% 		innertopmargin=5pt,
% 		innerbottommargin=8pt,
% 		innerleftmargin=8pt,
% 		leftmargin=-2pt,
% 		skipbelow=2pt,
% 		nobreak
% 	},
% 	headfont=\normalfont\bfseries\color{huidarkblue}
% ]{huiblueblox}
% \declaretheorem[style=huiblueblox,name=Example,within=chapter]{example}

% \declaretheoremstyle[
% 	mdframed={
% 		backgroundcolor=huilightblue,
% 		linecolor=huidarkblue,
% 		rightline=false,
% 		topline=false,
% 		bottomline=false,
% 		linewidth=2pt,
% 		innertopmargin=5pt,
% 		innerbottommargin=8pt,
% 		innerleftmargin=8pt,
% 		leftmargin=-2pt,
% 		skipbelow=2pt,
% 		nobreak
% 	},
% 	headfont=\normalfont\bfseries\color{huidarkblue}
% ]{huiblueblox}
\declaretheorem[style=huiblueblox,name=Problem,within=chapter]{prob}



% \declaretheoremstyle[
% 	mdframed={
% 		backgroundcolor=huilightblue,
% 		linecolor=huidarkblue,
% 		rightline=false,
% 		topline=false,
% 		bottomline=false,
% 		linewidth=2pt,
% 		innertopmargin=5pt,
% 		innerbottommargin=8pt,
% 		innerleftmargin=8pt,
% 		leftmargin=-2pt,
% 		skipbelow=2pt,
% 		nobreak
% 	},
% 	headfont=\normalfont\bfseries\color{huidarkblue}
% ]{huiblueblox}
\declaretheorem[style=huiblueblox,name=Exercise,within=chapter]{exer}
\declaretheorem[style=huipinkbox, name=Corollary, within=chapter]{cor}





















\newcommand{\nirwarnsymbol}{%
	\begin{tikzpicture}[baseline=(x.base)]
		\draw[rounded corners=.01em] (-.05em,-1.07em)rectangle(.05em,.78em);
		\draw[fill=white,rounded corners=1.3] (0,.75em)--(.75em,0)--(0,-.75em)--(-.75em,0)--cycle;
		\draw[line width=0.2mm, line cap=round](-.4em,-1.07em)--(.4em,-1.07em);
		\node(x) at (0,0em) {};
		% Thank you https://tex.stackexchange.com/a/262510
		\draw[
			line cap=but,
			line join=round,
			x=.5em,
			line width=0.5mm,
			y=1*(height("Z")-\pgflinewidth)*(1-sin(10)),
			rotate=-10,
			rounded corners=1.5pt,
		](-0.57, 0.57) -- (0.57, 0.57) -- (-0.57, -0.57) -- (0.57, -0.57);
	\end{tikzpicture}%
}

%%%%%%%%%%%%%%%%%%%%%%%%%%%%%%%%%%%%%%%%%%%% MARGINS
\usepackage{marginnote}
% Thank you https://tex.stackexchange.com/a/472882
% Makes marginnotes always appear on the left, apparently
%
\makeatletter
\long\def\@mn@@@marginnote[#1]#2[#3]{%
	\begingroup
		\ifmmode\mn@strut\let\@tempa\mn@vadjust\else
			\if@inlabel\leavevmode\fi
			\ifhmode\mn@strut\let\@tempa\mn@vadjust\else\let\@tempa\mn@vlap\fi
		\fi
		\@tempa{%
			\vbox to\z@{%
				\vss
				\@mn@margintest
				\if@reversemargin\if@tempswa
						\@tempswafalse
					\else
						\@tempswatrue
				\fi\fi

					\llap{%
						\vbox to\z@{\kern\marginnotevadjust\kern #3
							\vbox to\z@{%
								\hsize\marginparwidth
								\linewidth\hsize
								\kern-\parskip
								%\mn@parboxrestore
								\marginfont\raggedleftmarginnote\strut\hspace{\z@}%
								\ignorespaces#1\endgraf
								\vss
							}%
							\vss
						}%
						\if@mn@verbose
							\PackageInfo{marginnote}{xpos seems to be \@mn@currxpos}%
						\fi
						\begingroup
							\ifx\@mn@currxpos\relax\else\ifx\@mn@currpos\@empty\else
									\kern\@mn@currxpos
							\fi\fi
							\ifx\@mn@currpage\relax
								\let\@mn@currpage\@ne
							\fi
							\if@twoside\ifodd\@mn@currpage\relax
									\kern-\oddsidemargin
								\else
									\kern-\evensidemargin
								\fi
							\else
								\kern-\oddsidemargin
							\fi
							\kern-1in
						\endgroup
						\kern\marginparsep
					}%
			}%
		}%
	\endgroup
}
\makeatother
%
% Mostly for todonotes
\renewcommand{\marginpar}{\marginnote}
%%%%%%%%%%%%%%%%%%%%%%%%%%%%%%%%%%%%%%%%%%%% /MARGINS

\definecolor{nirlightred}{RGB}{250, 220, 220}
\definecolor{nirdarkred}{HTML}{f40000}
\declaretheoremstyle[
	mdframed={
		backgroundcolor=nirlightred,
		linecolor=nirdarkred,
		rightline=false,
		topline=false,
		bottomline=false,
		linewidth=2pt,
		innertopmargin=5pt,
		innerbottommargin=8pt,
		innerleftmargin=8pt,
		leftmargin=-2pt,
		skipbelow=2pt,
		nobreak
	},
	headfont=\normalfont\bfseries\color{nirdarkred}
]{nirredbox}

% \makeatletter
% \declaretheorem[
% 	style=nirredbox,
% 	name=Warning,
% 	sibling=thm,
% 	% without \leavevmode, the first item in a list gets misformatted
% 	postheadhook={\leavevmode\marginnote{\nirwarnsymbol}[-3pt]%
% 	\ifthmt@thisistheone% restatable makes alignment weird
% 		\hspace{-2.2pt}%
% 	\fi}
% ]{warn}
% \makeatother

\newcommand{\nirideasymbol}{%
	\begin{tikzpicture}[baseline=(x.base)]
		\draw[rounded corners=.01em] (-.05em,-1.07em)rectangle(.05em,.78em);
		\draw[fill=white,rounded corners=1.3] (0,.75em)--(.75em,0)--(0,-.75em)--(-.75em,0)--cycle;
		\draw[line width=0.2mm, line cap=round](-.4em,-1.07em)--(.4em,-1.07em);
		\node(x) at (0,0em) {};
		\node at (0,0em) {{\textbf{!}}};
	\end{tikzpicture}%
}
\renewcommand{\nirwarnsymbol}{%
	\begin{tikzpicture}[baseline=(x.base)]
		\draw[rounded corners=.01em] (-.05em,-1.07em)rectangle(.05em,.78em);
		\draw[fill=white,rounded corners=1.3] (0,.75em)--(.75em,0)--(0,-.75em)--(-.75em,0)--cycle;
		\draw[line width=0.2mm, line cap=round](-.4em,-1.07em)--(.4em,-1.07em);
		\node(x) at (0,0em) {};
		% Thank you https://tex.stackexchange.com/a/262510
		\draw[
			line cap=but,
			line join=round,
			x=.5em,
			line width=0.5mm,
			y=1*(height("Z")-\pgflinewidth)*(1-sin(10)),
			rotate=-10,
			rounded corners=1.5pt,
		](-0.57, 0.57) -- (0.57, 0.57) -- (-0.57, -0.57) -- (0.57, -0.57);
	\end{tikzpicture}%
}
\makeatletter
\declaretheorem[
	style=nirredbox,
	name=Idea,
	sibling=thm,
	% without \leavevmode, the first item in a list gets misformatted
	postheadhook={\leavevmode\marginnote{\nirideasymbol}[-3pt]%
	\ifthmt@thisistheone% restatable makes alignment weird
		\hspace{-2.2pt}%
	\fi}
]{idea}

\declaretheorem[
	style=nirredbox,
	name=Warning,
	sibling=thm,
	% without \leavevmode, the first item in a list gets misformatted
	postheadhook={\leavevmode\marginnote{\nirwarnsymbol}[-3pt]%
	\ifthmt@thisistheone% restatable makes alignment weird
		\hspace{-2.2pt}%
	\fi}
]{warn}
\makeatother

\title{Calc III Sections
\\ 
\vspace{0.4cm}
\large Fall 2025}




\date{\today}
\author{Hui Sun}


\begin{document}

\maketitle

% \tableofcontents
\newpage


\begin{center}
    \Large Calc III-Week 4 (9/15-9/19)
\end{center}

\renewcommand\thesection{\arabic{section}}

\noindent
Topics: (1) Partial derivatives. (2) Definition of total derivatives.

\begin{prob}
    Compute $\frac{\partial f}{\partial x}, \frac{\partial f}{\partial y}$ for the following functions:
    \begin{enumerate}
        \item \begin{equation*}
            x^3y^4-xy^2
        \end{equation*}
        \item \begin{equation*}
            x^2\sin(2y)+3
        \end{equation*}
        \item \begin{equation*}
            \ln\left(\frac{y}{x}\right)+\ln\left(\frac{1}{x+y}\right)-\ln\left(\frac{x}{2}\right)
        \end{equation*}
        You may use the following identities to simply the equation first:
        \begin{equation*}
            \ln\left(\frac{a}{b}\right)=\ln a-\ln b, \quad \ln\left(\frac{1}{a}\right)=-\ln a
        \end{equation*}
    \end{enumerate}
\end{prob}
% \begin{proof}
%     We have 
%     \begin{enumerate}
%         \item \begin{equation*}
%             \partial x: 3x^2y^4-y^2, \quad \partial y: 4x^3y^3-2xy
%         \end{equation*}
%         \item \begin{equation*}
%             \partial x: 2x\sin(2y), \quad \partial y: 2x^2\cos(2y)
%         \end{equation*}
%         \item \begin{equation*}
%             \partial x: -\frac{2}{x}-\frac{1}{x+y}, \quad \partial y: \frac{1}{y}-\frac{1}{x+y}
%         \end{equation*}
%     \end{enumerate}
% \end{proof}




\begin{defn}[tangent plane]
    Let $f:\R^2\to\R$ be differentiable at $(x_0,y_0)$, then the \textbf{tangent plane} to the graph $f$ in $\R^3$ is the given by 
    \begin{equation*}
        z=f(x_0,y_0)+\frac{\partial f}{\partial x}\bigg\vert_{(x_0,y_0)}(x-x_0)+\frac{\partial f}{\partial y}\bigg\vert_{(x_0,y_0)}(y-y_0)
    \end{equation*}
\end{defn}

\begin{prob}
    Compute the plane tangent to the graph of $f(x,y)=x^2y+2xy-y^2$ at $(1,2)$.
\end{prob}
% \begin{proof}
%     We have 
%     \begin{equation*}
%         \frac{\partial f}{\partial x}(1,2)=2xy+2y\vert_{(1,2)}=8, \quad \frac{\partial f}{\partial y}(1,2)=x^2+2x-2y\vert_{(1,2)}=-1
%     \end{equation*}
%     and $f(1,2)=2$, thus the plane is given by
%     \begin{equation*}
%         z=2+8(x-1)-(y-2)
%     \end{equation*}
%     i.e., $z=8x-y-4$.
% \end{proof}

\begin{defn}[derivative for two variables]
    Let $f:\R^2\to\R$, then $f$ is said to be \textbf{differentiable} at $(x_0,y_0)$ if $\frac{\partial f}{\partial x}, \frac{\partial f}{\partial y}$ exist at $(x_0,y_0)$ and if 
    \begin{equation*}
        \lim_{(x,y)\to(x_0,y_0)}\frac{f(x,y)-\mathcal{P}(x,y)}{\|(x,y)-(x_0,y_0)\|}=0
    \end{equation*}
    where $\mathcal{P}(x,y)=f(x_0,y_0)+\frac{\partial f}{\partial x}\bigg\vert_{(x_0,y_0)}(x-x_0)+\frac{\partial f}{\partial y}\bigg\vert_{(x_0,y_0)}(y-y_0)$ is the tangent plane to $f$ at $(x_0,y_0)$.
\end{defn}


\begin{defn}[derivative for $n$ variables]
    Let $f:\R^n\to\R$, then the \textbf{gradient} of $f$, denoted as $\nabla f$ is given by 
    \begin{equation*}
        \nabla f=\begin{bmatrix}
            \frac{\partial f}{\partial x_1}, \dots, \frac{\partial f}{\partial x_n}
        \end{bmatrix}
    \end{equation*}
    is a $1\times n$ matrix. And $f$ is said to be \textbf{differentiable} at $x_0\in\R^n$ if 
    \begin{equation*}
        \lim_{x\to x_0}\frac{|f(x)-f(x_0)-\nabla f(x_0)(x-x_0)|}{\|x-x_0\|}=0
    \end{equation*}
    and the derivative of $f$ is exactly the gradient $\nabla f$ at $x_0$.
\end{defn}

\begin{defn}[derivative for $m$ outputs]
    Let $f:\R\to\R^m$, where $f(x)=(f_1(x),\dots, f_m(x))$, then let $T$ denote the $n\times 1$ matrix 
    \begin{equation*}
        T=\begin{bmatrix}
            \frac{df_1}{dx}(x_0)\\
            \frac{df_2}{dx}(x_0)\\
            \dots\\
            \frac{df_m}{dx}(x_0)
        \end{bmatrix}
    \end{equation*}
    Then $f$ is said to be \textbf{differentiable} at $x_0$ if 
    \begin{equation*}
        \lim_{x\to x_0}\frac{|f(x)-f(x_0)-T(x-x_0)|}{|x-x_0|}=0
    \end{equation*}
    and the matrix $T$ is the derivative at $x_0$.
\end{defn}
\begin{example}
    Let $f(x)=(x^2, 2x, -x)$, then 
    \begin{equation*}
        T=Df(1)=\begin{bmatrix}
            2x\\
            2\\
            -1
        \end{bmatrix}=
        \begin{bmatrix}
            2\\
            2\\
            -1
        \end{bmatrix}
    \end{equation*}
\end{example}


\begin{defn}[derivative for general functions]
    Let $f: \R^n\to\R^m$, let $T$ be the $m\times n$ matrix with entries $\partial f_i/\partial x_j$ evaluated at $x_0\in \R^n$. Then  $f$ is said to be \textbf{differentiable} at $x_0$ if 
    \begin{equation*}
        \lim_{x\to x_0}\frac{\|f(x)-f(x_0)-T(x-x_0)\|}{\|x-x_0\|}=0
    \end{equation*}
    then $f$ is differentiable at $x_0$, and the matrix $T$ is the derivative at $x_0$. Note that $T$ loosk like 
    \[
T=Df(x_0) , \quad Df=
\begin{bmatrix}
\frac{\partial f_1}{\partial x_1} & \cdots & \frac{\partial f_1}{\partial x_n} \\
\vdots & \ddots & \vdots \\
\frac{\partial f_m}{\partial x_1} & \cdots & \frac{\partial f_m}{\partial x_n}
\end{bmatrix}
\]
\end{defn}
\begin{example}
    Let $f(x,y,z)=(ze^x, -ye^z)$, then 
    \[
Df(x, y, z) = 
\begin{bmatrix}
ze^x & 0 & e^x \\
0 & -e^z & -ye^z
\end{bmatrix}
\]
\end{example}


% \begin{warn}
%     It is not enough to just know the partial derivatives exist at $(x_0,y_0)$. Consider the following example: $f(x,y)=\sqrt{|xy|}$, note that $f$ is not differentiable at $(0,0)$.
% \end{warn}
% \begin{proof}
%     The partial derivatives are 
% \end{proof}







\begin{prob}
    Let \( f : \mathbb{R}^3 \to \mathbb{R} \) be defined by  
\[
f(x, y, z) = x^2 y + y \sin(z) + z e^x.
\]  
Compute the gradient of $f$ at $(1,2,0)$.
\end{prob}
% \begin{proof}
%     You can compute the partial derivatives 
%     \begin{align*}
%         \frac{\partial f}{\partial x} &= 2xy + ze^x, \\
%         \frac{\partial f}{\partial y} &= x^2 + \sin(z), \\
%         \frac{\partial f}{\partial z} &= y \cos(z) + e^x.
%         \end{align*}
%     thus the gradient is 
%     \[
% \nabla f = \left( \frac{\partial f}{\partial x}, \frac{\partial f}{\partial y}, \frac{\partial f}{\partial z} \right)=(2xy+ze^x, x^2+\sin z, y\cos z+e^x)
% \]
% Hence 
% \begin{equation*}
%     \nabla f(1,2,0)=(4, 1, 2+e)
% \end{equation*}

% \end{proof}


















\end{document}